\documentclass[../Vibro-Acoustic Signatures and Detection Algorithms for Stored Product Insects.tex]{subfiles}
\begin{document}

\chapter{Methodology}

\section{Acoustic Stored Product Insect Detection System (A-SPIDS)}

The Acoustic Stored Product Insect Detection System (A-SPIDS) was designed as a low-cost, portable, fast, and non-destructive alternative method to inspect stored products for insect infestations. The system is made up of three components. The main component is a box-in-a-box sampling container housed within a chamber outfitted with insulation, data acquisition and transmission electronics, and an integrated battery. A removable, acoustically treated lid, allows for easy access to the sampling container and can be securely closed to create a soundproof seal. The main component body can be set on a base that provides wireless charging and further vibration isolation. The external shell of the full system is made of 3D printed acrylonitrile butadiene styrene (ABS) plastic. The total system weighs \SI{9}{\kilo\gram} and has a length, height, and width of \SI{30}{\centi\meter}, \SI{24}{\centi\meter}, and \SI{19}{\centi\meter}, respectively. Figure \ref{fig:aspids_cad} shows an annotated rendered CAD model of A-SPIDS:

\begin{figure}[!htbp]
	\centering
	\includegraphics[width=\textwidth]{figures/methods/annotated_cad.pdf}
	\caption{Annotated rendered CAD model of A-SPIDS \label{fig:aspids_cad}}
\end{figure}

The integrated receptacle \SI{26}{\centi\meter} long, \SI{15}{\centi\meter} wide, and \SI{3.5}{\centi\meter} tall, capable of holding \SIrange{1}{1.5}{\litre} of product. The bottom is lined with piezoelectric discs to capture the mechanical sounds and vibrations generated by the insects moving, masticating, and mating within the product. This container size not only satisfies inspection sampling requirements but also stays within range for optimal insect detection based on the success of previous piezoelectric based sensors to detect \textit{Sitophilus oryzae} and \textit{Tribolium castaneum} in grain up to \SIrange{10}{15}{\centi\meter} and \SI{18.5}{\centi\meter} away, respectively, from the sensor \cite{vickSoundDetectionStoredProduct1988,hagstrumImprovingStoredProduct2019}. The maximum distance from the piezoelectric discs to any portion of the sample is within the range even when insects are active only on the substrate surface \cite{hagstrumEcologyIPMInsects2010}.

% The bottom of this receptacle is lined with 32 piezoelectric discs arranged in four groups of eight. This container size not only satisfies inspection sampling requirements but also stays within range for optimal insect detection based on the success of previous piezoelectric based sensors to detect the rice weevil, \textit{S. oryzae}, and the adult red flour beetle, \textit{T. castaneum}, in grain up to \SIrange{10}{15}{\centi\meter} and \SI{18.5}{\centi\meter} away, respectively, from the sensor \cite{vickSoundDetectionStoredProduct1988,hagstrumImprovingStoredProduct2019}. The maximum distance from the piezoelectric discs to any portion of the sample is set within the range even if the insects are active only on the surface of the substrate \cite{hagstrumEcologyIPMInsects2010}.

\newpage
The sides of the container and the bottom below the sensor array are insulated with a layer of \SI{5}{\milli\meter} of acoustic foam \cite{QuietBarrierHD2020} and \SI{5}{\milli\meter} layer of mass-loaded vinyl (MLV) \cite{QuietBarrierSpecialty2026}. Figure \ref{fig:aspids_insulation_diagram} depicts the cross-section of the container and the placement of the insulation layers.

\begin{figure}[!htbp]
	\centering
	\includegraphics[width=\textwidth]{figures/methods/insulation.pdf}
	\caption{Diagram of A-SPIDS insulation layout \label{fig:aspids_insulation_diagram}}
\end{figure}

The discs (PUI Audio, model AB2746B) are \SI{27}{\milli\meter} in diameter, \SI{0.51}{\milli\meter} thick, and made of piezoelectric ceramic lead zirconate titanate and brass. They have a resonant frequency of \SI{4600(500)}{\hertz}, a capacitance of \SI{16}{\nano\farad} $\pm$ \SI{30}{\percent} at \SI{1000}{\hertz}, a maximum resonant impedance of \SI{250}{\ohm}, and a maximum input voltage rating of \SI{20}{\volt}p-p \cite{BendersAB2746B}. The 32 discs are arranged in four groups of eight, separated by \SI{65}{\milli\meter}, corresponding to channels zero through three based on their respective location in the platform from left to right.

% The discs are \SI{27}{\milli\meter} in diameter, \SI{0.33}{\milli\meter} thick, and made of piezoelectric ceramic lead zirconate titanate and brass \cite{}. The four groups correspond to channels zero through three in the recordings, complementary to their respective location in the system from left to right seperated by \SI{65}{\milli\meter}. The piezoelectric discs capture the mechanical sounds and vibrations generated by the insects moving, masticating, and mating within the product. 

Traditional acoustic microphones (PUI Audio, model POM-2735P-R) are placed throughout the body of the system to capture external noise and vibration. The microphones have a diameter of \SI{6}{\milli\meter}, a height of \SI{4.7}{\milli\meter}, and a sensitivity of \SI{-35(2)}{\decibel} \cite{puiaudioPOM2735PRMicrophone2024}. Microphone channels four and five are colocated with the piezoelectric channels zero and one, respectively, while channel six is placed within the system body, and channel seven is mounted on the exterior of the system. Figure \ref{fig:aspids_sensor_diagram} shows the layout of the sensors.

\begin{figure}[!htbp]
	\centering
	\includegraphics[width=\textwidth]{figures/methods/sensor_module_layout.pdf}
	\caption{Diagram of A-SPIDS sensor layout \label{fig:aspids_sensor_diagram}}
\end{figure}

\newpage
The main body also contains the supportive hardware for data acquisition, communication, and power. The piezoelectric discs and microphones output to a proprietary amplifier and then converted to an I2S stream using a Texas Instruments PCM1808 analog-to-digital converter (ADC) \cite{PCM1808TexasInstruments}. The data is collected by a Raspberry Pi 3 A+ \cite{industriesRaspberryPiModel} and the I2S stream is digitized using a miniDSP MCHStreamer digital USB sound card \cite{MCHStreamerUSBAudio} with a \SI{24}{\kilo\hertz} sampling rate. The Raspberry Pi transmits data wirelessly to an acquisition computer through a Mikrotik mAP model RBmAP2nD \cite{MAPMikroTik}. The electronics are arranged on an electro-magnetic interference (EMI) shielding plate. Figure \ref{fig:aspids_elecronics_diagram} diagrams the electronics in A-SPIDS.

\begin{figure}[!htbp]
	\centering
	\includegraphics[width=\textwidth]{figures/methods/electronics_layout.pdf}
	\caption{Diagram of A-SPIDS electronics \label{fig:aspids_elecronics_diagram}}
\end{figure}

An integrated battery can maintain the system for several hours and recharges using a receiver coil when positioned on the platform base which is powered by a \SI{24}{\volt} power supply. The sensor can be lifted off and operated without the charging platform.

% During operation, the inspector places the sample material into the system enclosure and secures the lid. Unless the commodity is already packaged and portioned to \SIrange{1}{1.5}{\liter}, the product can be dispensed into a plastic bag. The inspector then initiates the data collection process until the system responds with the detection results. 

Figure \ref{fig:aspids_photos} shows photographs of A-SPIDS in the closed and open configurations.

\begin{figure}[!htbp]
	\centering
	\begin{subfigure}{0.45\textwidth}
		\centering
		\includegraphics[width=\textwidth]{figures/methods/aspids1.jpg}
		\caption{Closed configuration}
		\label{fig:aspids_open}
	\end{subfigure}%
	\qquad
	\begin{subfigure}{0.45\textwidth}
		\centering
		\includegraphics[width=\textwidth]{figures/methods/aspids2.jpg}
		\caption{Open configuration}
		\label{fig:aspids_closed}
	\end{subfigure}%
	\caption{A-SPIDS in (a) closed and (b) open configurations}%
	\label{fig:aspids_photos}
\end{figure}

\newpage
\section{Experimental Materials and Data Acquisition}

\subsection{Tested Materials}

A variety of stored product materials were tested with and without insect specimen to evaluate the performance of A-SPIDS across common commodities that ranged in processing level, density, and particle size. Hulled oatmeal, \textit{Avena sativa}, short-grain white rice, \textit{Oryza sativa}, wheat, \textit{Triticum aestivum}, in the form of processed all-purpose flour and unprocessed hard red wheat groats, and corn, \textit{Zea mays}, as the processed cereal Corn Flakes, were purchased from local grocery stores. Table \ref{tab:materials} provides a summary of the tested materials including their density, particle size, and hardness value. The density values were calculated by measuring the mass of \SI{1}{\liter} of each material. The particle, or kernel, size and the hardness values, the resistance of the material to deformation, were obtained from literature. 

\begin{table}[!htbp]
	% \small
	\caption{Description of tested materials \label{tab:materials}}
	\centering
	\small
	\begin{tabular}{@{}lrrrl@{}}
		\toprule
		\textbf{Common Name} & \textbf{Density [$\unit{\kilo\gram\per\cubic\meter}$]} & \textbf{Particle Size}                                              & \textbf{Hardness Value [$\unit{\kilo\gram f \per\milli\meter\cubed}$]}                                                                                                                   \\
		\midrule
		Oatmeal              & 338.14                                                 & \qtyrange[range-units=single,range-phrase=-]{6}{10}{\milli\meter}   & 10-14                                                                  & \cite{gatesMechanicalPropertiesOats2008}                                                                        \\
		Rice                 & 770.00                                                 & \qtyrange[range-units=single,range-phrase=-]{4}{5.5}{\milli\meter}  & 15-20                                                                  & \cite{nagatoGrainTextureRice1963}                                                                               \\
		Flour                & 600.00                                                 & \qtyrange[range-units=single,range-phrase=-]{30}{150}{\micro\meter} & 1                                                                      & \cite{glennMechanicalPropertiesStarch1992} \\
		Wheat Groats         & 790.00                                                 & \qtyrange[range-units=single,range-phrase=-]{5}{8}{\milli\meter}    & 18-25                                                                  & \cite{greerWaterAbsorptionWheat1959, morrisCollaborativeAnalysisWheat2011} \\
		Corn Flakes          & 118.35                                                 & \qtyrange[range-units=single,range-phrase=-]{5}{20}{\milli\meter}   & 8-12                                                                   & \cite{chaunierPhysicalSensoryEvaluation2005, lewickiInfluenceWaterActivity2007, ahnStudyOptimalThickness2020}   \\
		\bottomrule
	\end{tabular}
\end{table}


\subsection{Live Insect Specimens}

Although stored products insects proliferate across a wide range of commodities, they are difficult to obtain for research purposes. The cowpea beetle, \textit{Callosobruchus maculatus}, the confused flour beetle, \textit{Tribolium confusum}, and the mealworm and darkling beetle, \textit{Tenebrio molitor}, are commercially available and were chosen based on their similarities to the sizes, qualities, and behavior to the species targeted during inspection. The insects were obtained from Rainbow Mealworms (Compton, CA) \cite{RainbowMealwormsCrickets}. The insects were reared in the laboratory at ambient temperatures and humidity conditions. The \textit{C. maculatus} and \textit{T. confusum} were raised in their respective preferred substrates of cowpea and wheat flour, respectively, while the \textit{T. molitor} were fed oatmeal. Table \ref{tab:insects} provides a summary of the average size, weight, and density values measured for the tested insect species.

% make the font smaller for the table
\begin{table}[!htbp]
	% \small
	\caption{Description of tested live specimens \label{tab:insects}}
	\centering
	\small
	\begin{tabular}{@{}llrrr@{}}
		\toprule
		\textbf{Scientific Name}  & \textbf{Common Name}  & \textbf{Size [\si{\milli\meter}]} & \textbf{Weight [\si{\milli\gram}]} & \textbf{Density [\si{\kilo\gram\per\liter}]} \\
		\midrule
		\textit{C. maculatus}     & Cowpea beetle         & \(4 \times 1 \times 1\)           & 2.1                                & 0.525                                        \\
		\textit{T. confusum}      & Confused flour beetle & \(4 \times 2 \times 2\)           & 1.0                                & 0.625                                        \\
		\textit{T. molitor}       & Darkling beetle       & \(15 \times 5 \times 3\)          & 160.0                              & 0.711                                        \\
		\textit{T. molitor} larva & Mealworm              & \(20 \times 2 \times 2\)          & 130.0                              & 1.625                                        \\
		\bottomrule
	\end{tabular}
\end{table}

The following subsections provide a detailed description of the biology, behavior, previous detection studies, and any acoustic characteristics of the tested insect species.

\newpage

\subsubsection{\textit{Callosobruchus maculatus}}

The cowpea weevil or cowpea seed beetle and also known as the pulse beetle, \textit{Callosobruchus maculatus}, is of the order Coleoptera, Chrysomelidae family, and Bruchnae subfamily. These leaf beetles feed and develop exclusively on the seed of legumes but commonly infests, as their common name suggests, the cowpea, \textit{Vigna unguiculata}, also known as the southern pea or black-eyed pea \cite{beckHandbookBeanBeetles2014}.

The \textit{C. maculatus} develop through holometabolous metamorphosis with egg, larval, pupal, and adult stages. Their life cycle can be range from \SIrange{25}{37}{\day} at \SI{30}{\celsius} and \SI{70}{\percent} relative humidity but depend on temperature and humidity conditions \cite{katooCallosobruchusMaculatusBiotic2025}.

Adults are small, round beetles, brownish with black, pale spots and stripes, and range from \SIrange{2}{4}{\milli\meter} in length \cite{tranConsequencesInbreedingCowpea1995, katooCallosobruchusMaculatusBiotic2025}. A distinct feature is their 11 sectioned antennas. The insects are sexual dimorphic, where adult males are slightly smaller than their female counterparts \cite{katooCallosobruchusMaculatusBiotic2025}. Although the \textit{C. maculatus} have wings, only some forms of the species are capable of flight enabling them to disperse to new locations, including directly to cowpea fields. Adults do not require food or water and spend their one to two week lifespan mating and laying eggs \cite{beckHandbookBeanBeetles2014}. Figure \ref{fig:cb_drawing} shows the dorsal, ventral, and lateral views of \textit{C. maculatus}.

\begin{figure}[!htbp]
	\centering
	\begin{subfigure}{0.25\textwidth}
		\centering
		\includegraphics[width=\textwidth]{figures/insect_pictures/maculatus.jpeg}
		\caption{Dorsal}
		\label{fig:cb_front}
	\end{subfigure}
	\hfill
	\begin{subfigure}{0.25\textwidth}
		\centering
		\includegraphics[width=\textwidth]{figures/insect_pictures/maculatus_ventral.jpeg}
		\caption{Ventral}
		\label{fig:cb_ventral}
	\end{subfigure}
	\hfill
	\begin{subfigure}{0.25\textwidth}
		\centering
		\includegraphics[width=\textwidth]{figures/insect_pictures/maculatus side.jpeg}
		\caption{Lateral}
		\label{fig:cb_side}
	\end{subfigure}

	\caption{Dorsal, ventral, and lateral views of \textit{Callosobruchus maculatus}}%
	\label{fig:cb_drawing}
\end{figure}

Females lay eggs on the surface of the seeds. A single female can oviposit 120 to 132 eggs during its lifetime. The eggs are small, \SI{0.6}{\milli\meter}, oval and transform from transparent to translucent yellowish-white over time \cite{katooCallosobruchusMaculatusBiotic2025}.

The larvae hatch out of the egg in around five days and immediately dig into the seed, leaving the empty or mottled feces filled egg casing on the surface. The embryonic stage takes place within the seed while the larvae molts and feeds on the cotyledons and embryonic tissues of the seed. During this pupation stage, the larvae is exarate until after forming into an adult, after four to five days, it perforates out of the seed through a visible \SIrange{1}{2}{\milli\meter} circular hole \cite{beckHandbookBeanBeetles2014,katooCallosobruchusMaculatusBiotic2025}. This boring and consumption behavior is similar to other internal feeders, such as the lesser grain borer, \textit{R. dominica}, and the rice weevil, \textit{S. oryzae}, that are targeted as problematic by the USDA.

% Upon hatching in around five days, larvae dig into the seed and feed on the cotyledon, the embriotic leaf that is used to feed the developing plant, as they develop until they bore out of the bean as adults within 25-30 days \cite{tranConsequencesInbreedingCowpea1995}. This boring behavior is similar to other internal feeders, such as the lesser grain borer, \textit{R. dominica}, and the rice weevil, \textit{S. oryzae}, that are targeted as problematic by the USDA. Flying versions of the species can disperse to new locations, including directly to cowpea fields. 

The rapid life cycle of the \textit{C. maculatus} enable the colonies reproduce until their food supply is depleted \cite{tranConsequencesInbreedingCowpea1995}. The species can cause up to \SI{90}{\percent} damage in three to six months \cite{agourBrocchiaCinereaDelile2022}. Even a \SIrange{1}{2}{\percent} infestation can result in \SI{80}{\percent} loss of product after six to eight months of storage. The cowpea is a major protein source of over 200 million people in Sub-Saharan Africa \cite{bittnerMonitoringCowpeaBruchid2018}. Africa produces \SI{95}{\percent} of the global cowpea, with Nigeria outputting \SI{61}{\percent} of the \SIrange{1.4}{1.7}{\mega\tonne} annually \cite{bamaiyiEffectCallosobruchusMaculatus2006}. In 1986, Nigeria reported an annual loss of fifty million dollars worth of food to \textit{C. maculatus} \cite{smarttCowpeaResearchProduction1986}. In West Africa, up to \SI{100}{\percent} loss was reported after several months of unmitigated infestation \cite{mohamedn.sallamInsectDamagePostHarvest2013}.

Inspection for \textit{C. maculatus} was successful using X-ray and near-infrared hyperspectral imaging \cite{chelladuraiDetectionCallosobruchusMaculatus2014, divyanthIdentificationGreenGram2022} and GC-MS identification of sex pheromones \cite{kalirameshDetectionInfestationCallosobruchus2013}. Acoustic detection of the \textit{C. maculatus} has been used as a model organism to study methods for other internal feeders \cite{thanthiangaVibrationsMediatePrudent1987, devereauAutomaticRapidScreening2003, njorogeOxygenConsumptionAcoustic2018, bittnerMonitoringCowpeaBruchid2018}.

The \textit{C. maculatus} larvae were successfully detected by the Ultrasonic Insect Detector (UID) \cite{shadeDetectionHiddenInsect1990}. A wideband ultrasonic transducer determined that the feeding energy was contained to a frequency range from \SIrange{5}{75}{\kilo\hertz} but noted that the most prominent frequencies were between \SIrange{7.5}{12.5}{\kilo\hertz} despite the potential of external noise interference in operational settings \cite{klaassenCharacterizationUltrasonicSignals1988}.

The Purdue Insect Feeding Biomonitor (PIFB) detected the vibration generated by the larval feeding in cowpeas with a narrowband ultrasonic transducer set at \SI{40(2)}{\kilo\hertz}. By applying a low-pass filter at \SI{20}{\hertz}, the received signal is demodulated to only pulses generated by the insect caused events which are then enumerated with a Schmitt trigger threshold detector. The analysis showed that the number of events measured corresponded to the number of adults emerging from the seeds \cite{shadeDetectionHiddenInsect1990}.

A PCB 352C15 high-frequency, ceramic shear accelerometer, with a flat broadband \SI{+-3}{\decibel} response up to \SI{25}{\kilo\hertz}, was able to observe the larval masticating at a lower frequency range of \SIrange{10}{35}{\kilo\hertz} when using a Z-score threshold set four standard deviations above the normal distribution of the noise to isolate impulse events \cite{bittnerMonitoringCowpeaBruchid2018}. These ultrasonic sensor studies were conducted in a laboratory where a single, confirmed infested, bean was placed directly on the sensor. The \textit{C. maculatus} produced pulses with amplitudes of \SI{84}{\decibel} and \SI{95}{\decibel} in chickpea and green gram, respectively, and durations of \SI{68}{\milli\second} when using CZN-15E omnidirectional microphones in a ply board, iron, and polyurethane foam \SI{280}{\milli\meter} by \SI{280}{\milli\meter} by \SI{480}{\milli\meter} chamber \cite{bangaBioacousticDetectionCallosobruchus2019}.

The bioacoustic characteristics of the \textit{C. maculatus} in red and green beans were measured using a probe sensor inserted into a \SI{10}{\kilo\gram} storage bin with a Noise Reduction Coefficient of 0.8. A high-pass filter was applied with a cutoff frequency of \SI{2000}{\hertz} and a volume threshold of \SI{-8}{\deci\bel} was used to isolate and enumerate the insect-generated pulses for statistical analysis. The species was found to be most active in feeding and movement sounds when the temperature was in between \SIrange{10}{38}{\degreeCelsius} and increased as it rose. The larvae frequency peak was at \SI{2875(95.11)}{\hertz} with a pulse duration of \SIrange{40}{60}{\milli\second} in red bean and \SI{3588.33(253.54)}{\hertz} with a pulse duration of \SIrange{50}{60}{\milli\second} in green bean \cite{changBioacousticDetectionFour2025}.

\newpage
\subsubsection{\textit{Tribolium confusum}}

The confused flour beetle, \textit{Tribolium confusum} Jacquelin du Val, is an arthropod in the order Coleoptera and family Tenebrionidae. The common name is derived from its similarity and misidentification with the red flour beetle, \textit{T. castaneum}, which is a more problematic pest of stored products \cite{baldwinConfusedFlourBeetle2005}. The species has African origins \cite{tyagiPostHarvestStoredProduct2019}, originally  habituating beneath tree bark and rotting wood, they became secondary colonizers of stored products with evidence of their infestations in ancient Egyptian flour urns dating 5000 years \cite{pointerTriboliumBeetlesModel2021}.

\textit{T. confusum} is a secondary pest that infests farinaceous, high starch, products high in vitamin B such as rice and flour but also feeds on dried plants and fruits, dairy products, seeds, nuts, animal feed, cotton, spices, and processed foods such as pasta, cake mix, and crackers \cite{fordResearchPopulationsTribolium1937,tyagiPostHarvestStoredProduct2019,kavallieratosBiologicalFeaturesPopulation2020}. It has become a global pest and is abundant in flour mills across the United States \cite{StoredGrainInsects2016}.

Adults are reddish-brown with clubbed antennas and range from \SIrange{2}{6}{\milli\meter} in length \cite{baldwinConfusedFlourBeetle2020,tyagiPostHarvestStoredProduct2019} and weighing \SIrange{148}{178}{\milli\gram} \cite{fordResearchPopulationsTribolium1937}. The average life of an adult is \SI{1}{year} \cite{tyagiPostHarvestStoredProduct2019}. They have fully developed wings, capable of flight, but are behaviorally reluctant to do so \cite{kavallieratosBiologicalFeaturesPopulation2020}. Males can be distinguished by the presence of a gland in the forelimbs \cite{pointerTriboliumBeetlesModel2021}. Figure \ref{fig:cf_drawing} shows the dorsal, ventral, and lateral views of \textit{T. confusum}.

\begin{figure}[!htbp]
	\centering
	\begin{subfigure}{0.25\textwidth}
		\centering
		\includegraphics[width=\textwidth]{figures/insect_pictures/confused flour dorsal.jpeg}
		\caption{Dorsal}
		\label{fig:cf_front}
	\end{subfigure}
	\hfill
	\begin{subfigure}{0.25\textwidth}
		\centering
		\includegraphics[width=\textwidth]{figures/insect_pictures/confused flour ventral.jpeg}
		\caption{Ventral}
		\label{fig:cf_ventral}
	\end{subfigure}
	\hfill
	\begin{subfigure}{0.25\textwidth}
		\centering
		\includegraphics[width=\textwidth]{figures/insect_pictures/confused flour lateral.jpeg}
		\caption{Lateral}
		\label{fig:cf_side}
	\end{subfigure}

	\caption{Dorsal, ventral, and lateral views of \textit{Tribolium confusum}}%
	\label{fig:cf_drawing}
\end{figure}

Females lay an average of 450 eggs which are covered by a sticky secretion that gets covered by flour and meal adhering to sides of sacks, boxes, and other containers \cite{tyagiPostHarvestStoredProduct2019,StoredGrainInsects2016}. Eggs develop in 3 days \cite{pointerTriboliumBeetlesModel2021}.

Larva develops in 20 days \cite{pointerTriboliumBeetlesModel2021} commonly through six to 11 instars \cite{fordResearchPopulationsTribolium1937}. The larvae are small, wiry, slender, cylindrical and wormlike. They grow to a length of \SIrange{4}{5}{\milli\meter} and are white with yellow tinges \cite{StoredGrainInsects2016}. Pupa develops in 4 days transforming from white, to yellow, until they become the adult \cite{pointerTriboliumBeetlesModel2021,StoredGrainInsects2016}.

% Adults are reddish-brown with clubbed antennas and range from \SIrange{2}{4}{\milli\meter} in length \cite{baldwinConfusedFlourBeetle2020} and weighing \SIrange{148}{178}{\milli\gram} \cite{fordResearchPopulationsTribolium1937}.

The species successful global distribution is attributed to its ability to breed in temperatures from \SIrange{19}{37.5}{\degreeCelsius} and at humidity levels below \SI{1}{\percent}. It also is highly tolerant to common insecticides \cite{kavallieratosBiologicalFeaturesPopulation2020}.

Infestations of \textit{T. confusum} often result in a grey tint, odor, production of toxic quinones, and encourage mold growth. The confused flour beetle eggs are white and microscopic, proving a challenge to detect among the material they settle. Their lifecycle takes 40 to 90 days and the adults are capable of living up to three years in moderate climate conditions \cite{mohamedn.sallamInsectDamagePostHarvest2013, baldwinConfusedFlourBeetle2020}. Their lifecycle can be prolonged in cooler temperatures \cite{StoredGrainInsects2016}.

The \textit{T. confusum} responds to visual stimuli where tall, dark silhouettes were effective for trapping in ambient lighting conditions along with olfactory-baits \cite{semeaoResponseTriboliumCastaneum2011}. The species was also monitored by measuring its respiration through gas chromatography \cite{carlsonRespirationMeasurementTribolium1966}. The Enzyme-Linked Immunochemical Assay (ELISA) technique with polyclonal antibodies was able to detect the presence of the \textit{T. confusum} larvae, adults, and feces in grain with a concentration from \SIrange{4}{68}{\micro\gram\per\milli\liter} \cite{krizkova-kudlikovaDevelopmentPolyclonalAntibodies2008}. The DNA of the species was used to develop species-specific primers for real-time PCR detection \cite{nowaczykMolecularTechniquesDetection2009}.

The most detailed bioacoustic surveying of \textit{T. confusum} larvae and adults were conducted using a piezoelectric disc integrated in the wall of a \SI{10}{\milli\meter} diameter and \SI{50}{\centi\meter} tall PVC vertical cylinder. The container was filled with wheat grain and five larvae, adults, and a combination of both were placed at distances of \SIrange{10}{30}{\centi\meter} from the sensor. The recordings were actively annotated to distinguish the recognizable noises of larva and adult insects eating the product and adults moving in the substrate \cite{mousaviAcousticDetectionPossibility2017}. 

Larval feeding was observed in the range of \SIrange{1700}{2700}{\hertz} with a peak at \SI{2400}{\hertz}. The frequency ranges for adults were observed from \SIrange{1300}{3000}{\hertz} with a peak of \SI{2000}{\hertz} during eating and \SIrange{1900}{3000}{\hertz} with a peak of \SI{2300}{\hertz} during movement. While the frequency response remained unchanged as the distance between the sensor and insect increased, the pulse duration and amplitude at the peak frequency decreased with distance. Even at a distance of \SI{30}{\centi\meter}, the system was capable of detecting the insect-generated noises with reasonable fidelity \cite{mousaviAcousticDetectionPossibility2017}.

At a \SI{30}{\centi\meter} distance, sole feeding larvae were observed to produce pulses from \SIrange{16}{21}{\milli\second} in duration and \SIrange{-38}{-35}{\deci\bel} in amplitude at the peak frequency. The pulses of the walking and feeding adults were observed to be \SIrange{10}{14}{\milli\second} and \SIrange{18}{24}{\milli\second} in duration, respectively, and \SIrange{-39}{-34}{\deci\bel} and \SIrange{-37}{-33}{\deci\bel} in amplitude at the peak frequency, respectively \cite{mousaviAcousticDetectionPossibility2017}.

During measurements of larvae and adults in combination, the duration, amplitude, and frequency range of the pulses remained consistent but with more observed activity. Although the system demonstrated a promising ability for deployment, its detection performance was not evaluated \cite{mousaviAcousticDetectionPossibility2017}.

A different detection system was implemented using piezoelectric sensor mounted at the tip of a \SI{14}{\centi\meter} probe that was inserted into glass jars with one, two, and \SI{10}{\kilo\gram} of wheat and infested with a one to two live insect specimen. After amplifying the signal through a portable acoustic emission amplifier, a threshold-based algorithm was developed using a Hilbert transform and statistical analysis of the insect pulses. The system was determined to be able to detect a single \textit{T. confusum} in \SI{1}{\kilo\gram} of wheat with an accuracy of \SI{73}{\percent} \cite{eliopoulosDetectionAdultBeetles2015}.

\newpage
\subsubsection{\textit{Tenebrio molitor}}

The \textit{Tenebrio molitor} Linnaeus, part of the order Coleoptera and family Tenebrionidae, is known in its larval form as the yellow mealworm and in its adult form as the darkling beetle, named after its nocturnal habits and affinity for dark and moist environments.

The adult darkling beetle can grow up to \SI{2}{\centi\meter} in length and weigh up to \SI{100}{\milli\gram}. They have a hard, shiny brown or black exoskeleton, with a parallel-sided elytra, finely punctured thorax, and longitudinally striated wing covers \cite{adamaki-sotirakiStoredproductPestPromising2025}. Females can lay from 500  up to 6000 eggs during their lifecycle \cite{StoredGrainInsects2016}. The \textit{T. molitor} eggs are white and bean shaped, and just like the \textit{T. confusum}, are covered in a secretion that sticks to the substrate \cite{StoredGrainInsects2016}. The eggs hatch in \SIrange{3}{9}{\day} after deposition with a hatching rate ranging from \SIrange{65}{80}{\percent} \cite{adamaki-sotirakiStoredproductPestPromising2025}.

The larvae are \SIrange{1}{2}{\centi\meter} long, have a dark brown head, and a segmented yellow body with a hard exoskeleton. The larval stage can last from 57 days to up to two years under low temperatures during which it can experience from nine to 23 instars \cite{munoz-seijasCharacterizationAllLife2024}. The high protein and fat content of the larvae make them a popular feed for pets, livestock, fishing bait, and now are being considered as an alternative protein source for human consumption \cite{slowikHarmfulBeneficialSymbionts2023,munoz-seijasCharacterizationAllLife2024,yanSystematicReviewMicrobiological2025}. During the pupal stage, which lasts five to 48 days, the \textit{T. molitor} has a rigid exoskeleton and does not eat
\cite{munoz-seijasCharacterizationAllLife2024}.

Figure \ref{fig:tm_drawing} shows the dorsal, ventral, and lateral views of the \textit{T. molitor} and its larval form.

\begin{figure}[!htbp]
	\centering
	\begin{subfigure}{0.2\textwidth}
		\centering
		\includegraphics[width=\textwidth]{figures/insect_pictures/DarklingBeetle.jpeg}
		\caption{Dorsal}
		\label{fig:tm_front}
	\end{subfigure}
	\hfill
	\begin{subfigure}{0.2\textwidth}
		\centering
		\includegraphics[width=\textwidth]{figures/insect_pictures/mealworm bottom.jpeg}
		\caption{Ventral}
		\label{fig:tm_ventral}
	\end{subfigure}
	\hfill
	\begin{subfigure}{0.2\textwidth}
		\centering
		\includegraphics[width=\textwidth]{figures/insect_pictures/Mealworm side.jpeg}
		\caption{Lateral}
		\label{fig:tm_side}
	\end{subfigure}
	\hfill
	\begin{subfigure}{0.2\textwidth}
		\centering
		\includegraphics[width=\textwidth]{figures/insect_pictures/larva.jpg}
		\caption{Larval}
		\label{fig:mealworm}
	\end{subfigure}
	\caption{Dorsal, ventral, and lateral views of \textit{Tenebrio molitor} and \textit{T. molitor} larva}%
	\label{fig:tm_drawing}
\end{figure}

% Although the \textit{T. molitor} is not considered a significant pest by the USDA, unmitigated, the larval form can heavily deplete a supply \cite{StoredGrainInsects2016}. The species can feed on various vegetable and animal matter, it has a preference for cereal and grain products \cite{adamaki-sotirakiStoredproductPestPromising2025}. 

As a result of the increased adoption of the \textit{T. molitor} as a food source for human consumption, insights into the foodborne pathogens associated with the insect and the potential for contamination of the product have been a focus of research. The \textit{T. molitor} has been found to harbor a variety of pathogens such as \textit{Staphylococcus aureus} and spore-forming bacteria such as \textit{Bacillus cereus}, \textit{Clostridium perfringens}, and \textit{Clostridium botulinum} \cite{slowikHarmfulBeneficialSymbionts2023, yanSystematicReviewMicrobiological2025}.

When infesting a commodity, darkling beetles release benzoquinones that result in an unpleasant odor \cite{mohamedn.sallamInsectDamagePostHarvest2013, vigneronImmuneDefensesBeneficial2019}. The \textit{T. molitor} can also disperse the fungus \textit{Fusarium proliferatum}, a pathogen is capable of synthesizing mycotoxins beauvericin, fusaproliferin, fusarins, and moniliformin that can cause a variety of health issues in humans and animals \cite{guoDisseminationFusariumProliferatum2018}.

% PTA-ELISA was able to detect the \textit{T. molitor} larvae, adults, and feces in grain with a concentration from \SIrange{17}{32}{\micro\gram \per \milli\liter}, respectively. 

In its larval form, the \textit{T. molitor} generated \SI{40}{\deci\bel} impulses with peak frequencies from \SIrange{100}{1500}{\hertz} and \SIrange{4000}{15000}{\hertz} measured using a Br\"{u}el \& Kj\ae r model 2603 microphone in an anechoic chamber with an attenuation of \SI{60}{\deci\bel}
\cite{arnettMeasurementAdventitiousSounds1968}.

\newpage
\subsection{Testing Procedure and Dataset Description}

During testing, \SI{1}{\liter} of clean material was transferred into a plastic bag and placed into the sensor container. After several minutes of silence were recorded to ensure the grains were settled, a single insect was introduced into the material. Table \ref{tab:recording} summarizes the duration of recordings for each insect-material pairing. The insect-material pairings are not fully balanced as a result of varying life stages and availability of the insects, the need to test in clean and uncontaminated materials, and the time and budget restraints of the project.

\begin{table}[!htbp]
	\caption[Duration of recordings for insect-material pairings]{Duration of recordings in minutes for insect-material pairings with silence. The duration of recordings with external noise are accommodated with parenthesis. \label{tab:recording}}
	\centering
	\small
	\begin{tabular}{@{}l*{6}{r}@{}}
		\toprule
		                          & \textbf{\textit{A. sativa}} & \textbf{\textit{O. sativa}} & \multicolumn{2}{c}{\textbf{\textit{T. aestivum}}} & \textbf{\textit{Z. mays}} & \multirow{2}{*}{\textbf{Total}}              \\
		                          & Oatmeal                     & White Rice                  & Wheat Groats                                      & Flour                     & Corn Flakes                     &            \\
		\midrule
		\textit{C. maculatus}     & 12                          & 258 (12)                    & 13                                                & 20 (1)                    & 17                              & 320 (12)   \\
		\textit{T. confusum}      & 39                          & 274 (16)                    & 189                                               & 55 (1)                    & 16                              & 573 (34)   \\
		\textit{T. molitor}       & 68                          & 15                          & 9                                                 & 8                         & 13                              & 113 (4)    \\
		\textit{T. molitor} larva & 46 (4)                      & 268 (132)                   & 51 (2)                                            & 128                       & 43 (5)                          & 536 (143)  \\
		\midrule
		No insect                 & 4                           & 219 (64)                    & 142 (2)                                           & 24                        & 23                              & 412 (66)   \\
		\midrule
		Total                     & 169 (4)                     & 1034 (224)                  & 403 (4)                                           & 235 (2)                   & 112 (5)                         & 1953 (239) \\
		\bottomrule
	\end{tabular}
\end{table}


A curated subset of the Stored Product Insect Dataset was extracted for the analysis. The portion features \SI{1}{\minute} recordings of different insect-material pairings, in silence, when an insect was present and observed moving with at least five impulses in visualization figures. This resulted in a total of 1329 recordings, corresponding to \SI{1329}{\minute}, including five control recordings without insects but with ambient noise. Table \ref{tab:subset} details the number of \SI{1}{\minute} recordings for each insect-material pairing used for analysis:

\begin{table}[!htbp]
	\caption[Description of subset used for analysis]{Number of 1 min records for each insect-material pairing used for analysis. \label{tab:subset}}
	\centering
	\small
	\begin{tabular}{@{}l*{7}{r}@{}}
		\toprule
		                          & \textbf{\textit{A. sativa}} & \textbf{\textit{O. sativa}} & \multicolumn{2}{c}{\textbf{\textit{T. aestivum}}} & \textbf{\textit{Z. mays}} & \multirow{2}{*}{\textbf{Total}}        \\
		                          & Oatmeal                     & Rice                        & Wheat Groats                                      & Flour                     & Corn Flakes                     &      \\
		\midrule
		\textit{C. maculatus}     & 8                           & 201                         & 11                                                & 16                        & 16                              & 252  \\
		\textit{T. confusum}      & 33                          & 235                         & 177                                               & 42                        & 15                              & 502  \\
		\textit{T. molitor}       & 54                          & 11                          & 8                                                 & 8                         & 11                              & 92   \\
		\textit{T. molitor} larva & 40                          & 239                         & 41                                                & 123                       & 35                              & 478  \\
		\midrule
		No Insect                 &                             & 5                           &                                                   &                           &                                 & 5    \\
		\midrule
		Total                     & 135                         & 691                         & 237                                               & 189                       & 77                              & 1329 \\ \bottomrule
	\end{tabular}
\end{table}

\newpage
\subsection{Simulated External Noise}

The expected usage of the system would be within a variety of environments such as storage facilities, warehouses, and ports, where machinery, vehicles, and people can create significant background noise. The sound pressure levels can range from \SIrange{35}{50}{\decibel A}, common for an office, up to \SI{110}{\decibel A} such as a factory \cite{mankinPerspectivePromiseCentury2011,claudiaNoiseSourcesCharacterization2012}.

The system was exposed to varying levels of natural and simulated, ambient and artificial noise with and without live insect specimens to examine the efficacy of the sound insulation as well as the influence on the detection, false positive, and false negative performance. The ambient noise was a result of the surrounding laboratory environment including equipment, HVAC units, and personnel as well as vehicular and aircraft traffic. Artificial noise, such as pre-recorded sweeps, white noise, loud conversations, crowds, and factory noise were played through a loudspeaker system placed \SI{1}{\meter} from the system. The noise was played at increased steps from \SIrange{60}{110}{\decibel A} and the noise levels were measured using a calibrated sound level meter, the Extech 407730 Digital Sound Level Meter manufactured by Teledyne FLIR, placed coaxially \SI{10}{\centi\meter} from the external microphone channel seven \cite{EXTECH407730Digital}.

\begin{figure}[!htbp]%
	\centering
	\includegraphics[width=\textwidth]{figures/methods/external_noise_layout.pdf}
	\caption{Layout of the external noise measurement setup}%
	\label{fig:external_noise_measurement}%
\end{figure}

Spectrograms of the recordings are featured in Appendix \ref{app:noise-spectrograms} \cite{sonically_soundCrowdCafeteria2022,viertelnachvierFactory22014,inspectorjAmbienceMachineFactory2017,usnationalarchivesPresidentRonaldReagans2015}.

\newpage
\section{Stored Product Insect Database (SPIDB)}

The Stored Product Insect Dataset (SPID) is a repository of audio recordings collected by the Stored Product Insect Detection System. The dataset features the audio recordings as well as metadata describing the recording system, noise conditions, insect species, and material substrates. The full \SI{100}{\giga\byte} dataset can be downloaded from Kaggle \cite{danielkadyrovStoredProductInsecta}. An excerpt of the dataset file structure is shown in Listing \ref{lst:spid_directory}:

\begin{listing}[htbp]
	\caption{File organization of the Stored Product Insect Dataset \label{lst:spid_directory}}

	\begin{minted}{text}
\---data
¦   spi.db
¦  
+---aspids
¦   ¦   metadata.csv
¦   ¦   
¦   +---2023
¦   ¦   ¦   +---24
¦   ¦   ¦   ¦   +---2023_05_24_19_41_19.559_t00100
¦   ¦   ¦   ¦   ¦       2023_05_24_19_41_19.559_t00100_Ch0_00001.wav
¦   ¦   ¦   ¦   ¦       2023_05_24_19_41_19.559_t00100_Ch0_00002.wav
	\end{minted}
\end{listing}

The audio recordings are organized by date and the file names describe the recording date, time, a string descriptor, sensor channel, and recording iteration. The filename is structured as follows: \texttt{YYYY} is the four-digit year, \texttt{MM} is the two-digit month, \texttt{DD} is the two-digit day, \texttt{HH} is the two-digit hour in 24-hour format, \texttt{MM} is the two-digit minute, \texttt{SS.fff} is the seconds and milliseconds, \texttt{ssssss} is a string descriptor of the recording conditions, \texttt{Ch\#} is the sensor channel number, and \texttt{iiiii} is the five-digit recording iteration. The file name structure is detailed in Listing \ref{lst:filename}.

\begin{listing}[htbp]
	\caption{File organization of the Stored Product Insect Dataset \label{lst:filename}}
	\begin{minted}{text}
YYYY_MM_DD_HH_MM_SS.fff_ssssss_Ch#_iiiii.wav
	\end{minted}
\end{listing}

The Stored Product Insect Database (SPIDB) provides an interface to access the data collected from the A-SPIDS platform. Acoustic data requires a rich set of metadata to provide context. This could range from information about the recording system such as the sensor type and sensitivity, environmental conditions such as background noise levels and descriptions, as well as details about the featured subjects. Commonly, datasets, like BugBytes, provide acoustic data as a set of audio files along with text-based files annotating each recording. Even when data is provided, it is often decoupled and requires significant effort to match the audio.

\newpage
SPIDB extends the Sound Organization for Network Integration, Collection, and Collaboration (SONICDB) framework, which was developed to address the challenge of organizing and sharing acoustic data when large amounts of associated metadata are required for analysis and interpretation \cite{rochManagementAcousticMetadata2016}. SONICDB leverages relational databases to link audio files with their corresponding minutiae, enabling efficient data management and retrieval as well as simplified sharing across different users and computers. Although any database reader or programming language can be used to interact with a database, an object-oriented wrapper was developed in Python using industry standard packages and methods. This interface was used to directly write A-SPIDS data to the database as well as throughout the data processing and analysis pipeline. Currently, the database is a locally stored SQLite file, which can be easily version-controlled and distributed. The interface can also support server-based SQL database systems such as MySQL and PostgreSQL \cite{danielkadyrovDkadyrovSonicdbV0122025}.

Installing the SPIDB package can be performed through the Pip Installs Packages (PIP) Python package manager or by directly downloading from the GitHub repository. The dataset can be downloaded directly from Kaggle or using a script that is provided with the SPIDB package. Listing \ref{lst:spidb_install} describes the commands to install the SPIDB package and utilize package to download the dataset.

\begin{listing}[htbp]
	\caption{Installing SPIDB package and downloading the database \label{lst:spidb_install}}
	\begin{minted}{bash}
# Install the SPIDB package
pip install spidb 
# Create a directory 
mkdir spid 
# Download the SPIDB database into the spid directory
spidb download --output-dir spid 
	\end{minted}
\end{listing}

SONICDB initializes with several preliminary models. The \texttt{File} model represents the audio files stored in the database, detailing the file path, filename, extension, sample rate, duration, start and end time, as well as the sensor and channel the files were recorded with. The sensor and channel properties are directly tied to the \texttt{Sensor} and \texttt{Channel} models, respectively. The \texttt{Sensor} model categorizes the name, subname, manufacturer, sensor type, and number of channels of the sensor as well as a list of associated \texttt{Channel} and \texttt{File} entities. The \texttt{Event} model facilitates storing details of the logged events, including start and end times, description, noise levels, and the associated \texttt{Material} and \texttt{Subject} entities. The \texttt{Subject} model captures information about the subjects being recorded. SPIDB extended it to include bespoke attributes such as scientific name, common name, weight, height, length, width, volume, and density. SPIDB also added the \texttt{Material} model to describe the substrate in which the subjects are placed.

\newpage
Figure \ref{fig:spidb} shows the schema for the SPIDB:

\begin{figure}[!htbp]%
	\centering
	\includegraphics[width=\textwidth]{figures/methods/database_diagram.pdf}
	\caption{Schema of the Stored Product Insect Database (SPIDB)}%
	\label{fig:spidb}%
\end{figure}

SONICDB uses popular and standard Python packages ensuring standardized database management and coding conventions. SQLAlchemy allows defining the database schema using Python classes, making it easier to interact with the database using an object-oriented approach. Additionally, Alembic is used for database migrations, enabling version control for the database schema \cite{fredstamComparingDatabaseManagement}. The data is extracted as Pandas DataFrames and Numpy Arrays. Audio data is handled through Librosa and PyDub. Listing \ref{lst:spidb_connection} shows an example of connecting to the database using Python once the dataset is downloaded and the user is within the directory containing the database file, \texttt{spi.db}:

\begin{listing}[htbp]
	\caption{Example of connecting to the SPIDB using Python \label{lst:spidb_connection}}
	\begin{minted}{python}
from spidb import spidb 

db = spidb.Database("spi.db")
	\end{minted}
\end{listing}

\newpage
Once available, a user can interact with the database using the defined models and perform various operations such as querying, inserting, updating, and deleting records. Listing \ref{lst:spidb_query} shows an example of querying the database for all the tested materials and then retrieving all recordings where the material is associated with a specific insect species, \textit{C. maculatus}:

\begin{listing}[htbp]
	\caption{Example of querying the SPIDB for a specific insect species \label{lst:spidb_query}}
	\begin{minted}{python}
# Query for all tested materials
materials = db.Material.query.all() #[Material <id>, ...]

# Retrieve all events for a specific insect species
species = "Callosobruchus maculatus"
events = db.session.query(spidb.Event).filter(
    spidb.Event.subject.name == species
    ).filter(
        spidb.Event.material.id.in_([m.id for m in materials])
    ).all() # [Event <id>, ...]
	\end{minted}
\end{listing}

Although the audio can be retrieved by accessing the \texttt{File} entity, independently or through association to the \texttt{Event}, \texttt{Sensor}, or \texttt{Channel} entities, SONICDB provides an accessible method to retrieve the data if the start, end, and channel are known. Listing \ref{lst:sonicdb_retrieval} shows an example of how to retrieve audio data from SONICDB.

\begin{listing}[htbp]
	\caption{Example of retrieving audio data from SONICDB \label{lst:sonicdb_retrieval}}
	\begin{minted}{python}
# Retrieve the Event entity for which audio is to be extracted
event = db.session.query(spidb.Event).one() # <Event id>

# Retrieve audio from channel 2 for the duration of the event
audio = db.get_audio(
    event.start,
    event.end,
    sensor=event.sensor,
    channel_number=2
) # <Audio id>
	\end{minted}
\end{listing}

The code used to generate the data and figures in this report is available in the accompanying GitHub repository described in Appendix \ref{app:scripts}. The full documentation, including examples and scripts, for SONICDB and SPIDB, is available on GitHub \cite{danielkadyrovDkadyrovSonicdbV0122025,danielkadyrovDkadyrovSpidbSPIDB2024, danielkadyrovDkadyrovSPIDAcousticV100a2026}.

\newpage

\end{document}